\begin{thebibliography}{99}

\bibitem{acemogluRestrepo2018}
Acemoglu, D., and Restrepo, P. (2018).
Artificial Intelligence, Automation and Work.
\textit{NBER Working Paper} No. 24196.

\bibitem{acemogluAutor2011}
Acemoglu, D., and Autor, D. (2011).
Skills, Tasks and Technologies: Implications for Employment and Earnings.
In \textit{Handbook of Labor Economics}, Vol. 4B, 1043--1171.
Elsevier.

\bibitem{autor2015}
Autor, D. H. (2015).
Why Are There Still So Many Jobs? The History and Future of Workplace Automation.
\textit{Journal of Economic Perspectives}, 29(3), 3--30.

\bibitem{bresnahanTrajtenberg1995}
Bresnahan, T. F., and Trajtenberg, M. (1995).
General Purpose Technologies: Engines of Growth?
\textit{Journal of Econometrics}, 65(1), 83--108.

\bibitem{brynjolfssonHitt2000}
Brynjolfsson, E., and Hitt, L. M. (2000).
Beyond Computation: Information Technology, Organizational Transformation and Business Performance.
\textit{Journal of Economic Perspectives}, 14(4), 23--48.

\bibitem{brynjolfssonRockSyverson2019}
Brynjolfsson, E., Rock, D., and Syverson, C. (2019).
Artificial Intelligence and the Modern Productivity Paradox: A Clash of Expectations and Statistics.
In \textit{The Economics of Artificial Intelligence: An Agenda}.
University of Chicago Press.

\bibitem{cominMestieri2018}
Comin, D., and Mestieri, M. (2018).
If Technology Has Arrived Everywhere, Why Has Income Diverged?
\textit{American Economic Journal: Macroeconomics}, 10(3), 137--178.

\bibitem{aghionJonesJones2017}
Aghion, P., Jones, B. F., and Jones, C. I. (2017).
Artificial Intelligence and Economic Growth.
\textit{NBER Working Paper} No. 23928.

\bibitem{korinekStiglitz2021}
Korinek, A., and Stiglitz, J. E. (2021).
Artificial Intelligence, Globalization, and Strategies for Economic Development.
\textit{NBER Working Paper} No. 28453.

\bibitem{noyZhang2023}
Noy, S., and Zhang, W. (2023).
Experimental Evidence on the Productivity Effects of Generative Artificial Intelligence.
\textit{Science}, 381(6654), 187--192.

\bibitem{brynjolfssonLiRaymond2023}
Brynjolfsson, E., Li, D., and Raymond, L. R. (2023).
Generative AI at Work.
\textit{NBER Working Paper} No. 31161.

\bibitem{eloundouEtAl2023}
Eloundou, T., Manning, S., Mishkin, P., and Rock, D. (2023).
GPTs are GPTs: An Early Look at the Labor Market Impact Potential of Large Language Models.
\textit{arXiv preprint} arXiv:2303.10130.

\bibitem{bai2009}
Bai, J. (2009).
Panel Data Models with Interactive Fixed Effects.
\textit{Econometrica}, 77(4), 1229--1279.

\bibitem{callaway2021}
Callaway, B., and Sant'Anna, P. H. C. (2021).
Difference-in-Differences with Multiple Time Periods.
\textit{Journal of Econometrics}, 225(2), 200--230.

\bibitem{chernozhukov2018}
Chernozhukov, V., Chetverikov, D., Demirer, M., Duflo, E., Hansen, C., Newey, W., and Robins, J. (2018).
Double/Debiased Machine Learning for Treatment and Structural Parameters.
\textit{The Econometrics Journal}, 21(1), C1--C68.

\bibitem{griliches1998}
Griliches, Z. (1998).
\textit{R\&D and Productivity: The Econometric Evidence}.
University of Chicago Press.

\bibitem{goldfarbTucker2019}
Goldfarb, A., and Tucker, C. (2019).
Digital Economics.
\textit{Journal of Economic Literature}, 57(1), 3--43.

\bibitem{hallJones1999}
Hall, R. E., and Jones, C. I. (1999).
Why Do Some Countries Produce So Much More Output Per Worker Than Others?
\textit{Quarterly Journal of Economics}, 114(1), 83--116.

\bibitem{imf2024}
International Monetary Fund (2024).
Gen-AI: Artificial Intelligence and the Future of Work.
\textit{IMF Staff Discussion Note}.

\bibitem{oecd2023ai}
OECD (2023).
\textit{OECD AI Policy Observatory: Measuring AI Adoption and Diffusion}.
OECD Publishing.

\bibitem{jolliffeCadima2016}
Jolliffe, I. T., and Cadima, J. (2016).
Principal Component Analysis: A Review and Recent Developments.
\textit{Philosophical Transactions of the Royal Society A}, 374(2065), 20150202.

\bibitem{kaufmannKraayMastruzzi2010}
Kaufmann, D., Kraay, A., and Mastruzzi, M. (2010).
The Worldwide Governance Indicators: Methodology and Analytical Issues.
\textit{World Bank Policy Research Working Paper} No. 5430.

\bibitem{feenstraInklaarTimmer2015}
Feenstra, R. C., Inklaar, R., and Timmer, M. P. (2015).
The Next Generation of the Penn World Table.
\textit{American Economic Review}, 105(10), 3150--3182.

\bibitem{worldbankWDI2024}
World Bank (2024).
\textit{World Development Indicators Database}.
World Bank.

\bibitem{wooldridge2010}
Wooldridge, J. M. (2010).
\textit{Econometric Analysis of Cross Section and Panel Data} (2nd ed.).
MIT Press.

\bibitem{wooldridge2021}
Wooldridge, J. M. (2021).
\textit{Two-Way Fixed Effects, the Two-Way Mundlak Regression, and Difference-in-Differences Estimators}.
Working Paper.

\bibitem{worldbank2021wdr}
World Bank (2021).
\textit{World Development Report 2021: Data for Better Lives}.
World Bank.

\end{thebibliography}
